% ============================================================
% Chapter: Final Concept Evaluation and Recommendations
% Assessed MLOs: MLO-1, MLO-2, MLO-3
% Excellent-grade rubric criteria:
% - Concept design selection and evaluation (10%): the final concept is judged
%   against the accumulated technical evidence, and recommendations are specific,
%   defensible, and consistent with the preceding analysis.
% ============================================================

\section{Final Concept Evaluation and Recommendations}
\pagecheck{3--4}

% Rubric Checklist: Concept design selection and evaluation (10%) with emphasis on final feasibility judgement.
\instructionplaceholder{Integrated evaluation}{Use this chapter to close the loop on concept selection. Do not repeat earlier chapters section by section. Instead, synthesise the consequences of the metocean basis, concept studies, mooring behaviour, numerical modelling, and probabilistic assessment into one coherent feasibility judgement.}

\subsection{Impact of analysis on design choices}
\instructionplaceholder{Required evidence}{State explicitly how the numerical and probabilistic results affected the concept. Identify what changed in dimensions, detailing, system behaviour, safety margin, or design philosophy because of the analysis work.}

\subsection{Technical feasibility and limitations}
\instructionplaceholder{Required evidence}{Make a direct statement about technical feasibility. If the concept is feasible only under certain assumptions, list them. If major limitations remain unresolved, state them clearly and explain whether they are design issues, data issues, or modelling issues.}

\subsection{Recommendations and lessons learned}
\instructionplaceholder{Required evidence}{End with the next actions required to improve confidence in the concept. These should be engineering recommendations, not generic reflections. Add a short set of lessons learned only if they sharpen the technical recommendation or explain a key design decision.}
